\documentclass[11pt,a4paper]{article}
\ifdefined\pdfmapfile
  \pdfmapfile{+lm.map}
  \pdfmapfile{+cm.map}
  \pdfmapfile{+symbols.map}
  \pdfmapfile{+euler.map}
\fi
\usepackage[T1]{fontenc}
\usepackage{lmodern,microtype}
\usepackage[left=25mm,right=25mm,top=24mm,bottom=25mm,headheight=14pt]{geometry}
\usepackage{amsmath,amssymb,amsthm,booktabs,tabularx,array}
\usepackage{xcolor,fancyhdr,enumitem}
\definecolor{inkblue}{HTML}{174C60}
\usepackage[unicode,colorlinks=true,linkcolor=inkblue,citecolor=inkblue,urlcolor=inkblue]{hyperref}
\hypersetup{pdftitle={Prime-counting exceptions, Farey permutations, and zeros of a theta-kernel Mellin transform},pdfsubject={Research manuscript in analytic number theory and analysis}}
\pagestyle{fancy}\fancyhf{}
\fancyhead[L]{\small\itshape Prime-counting exceptions and theta-kernel zeros}
\fancyhead[R]{\small Research manuscript}\fancyfoot[C]{\thepage}
\renewcommand{\headrulewidth}{0.25pt}
\newtheorem{theorem}{Theorem}[section]
\newtheorem{proposition}[theorem]{Proposition}
\newtheorem{lemma}[theorem]{Lemma}
\newtheorem{corollary}[theorem]{Corollary}
\setlist{topsep=4pt,itemsep=2pt}
\setcounter{tocdepth}{1}
\setlength{\emergencystretch}{2em}
\raggedbottom
\begin{document}
\thispagestyle{plain}
\begin{center}
{\small\sffamily\color{inkblue} ANALYTIC NUMBER THEORY AND ANALYSIS}\\[9pt]
{\LARGE\bfseries Prime-counting exceptions, Farey permutations,}\\[4pt]
{\LARGE\bfseries and zeros of a theta-kernel Mellin transform}\\[10pt]
{\large Unconditional results and consequences of the Riemann hypothesis}\\[9pt]
{\small Research manuscript \quad $\cdot$ \quad 21 November 2026}
\end{center}

\begin{abstract}
Erd\H{o}s and Richards asked in 1977 whether, for a set of integers $n$
of natural density one, $\pi(n+h)-\pi(n)\le\pi(h)$ holds for every
integer $2\le h\le n$. We give an unconditional proof. For each fixed
$K\ge2$, lengths below the $K$th prime determine the set of exceptions
up to $O_K(x/(\log x)^{K+1})$ starts below $x$. Published optimal admissible diameters through size 342 then give $B(x)\ll x/(\log x)^{351}$ for the total exception count. A separate argument
answers Tom\'as Garc\'ia's conjecture on the discrepancy averaged over
permutations of prescribed gaps, with the sharp uniform lower constant.
Under the ordinary Riemann hypothesis the proportion of permutations of
the Farey gaps as regular as their natural ordering decays at least as
fast as $\exp(-c_\eta n^{1-\eta})$ for every $0<\eta<1$.
We also give a computer-assisted counterexample to Csordas's 1998
question about the zeros of a meromorphically continued Mellin transform.
The source contains the interval calculation and the bounds for the
omitted terms. Additional results concern a positive-definiteness
question of Csordas, the oscillation of Nicolas's totient sum over zeta
zeros, and primes between real powers with exponent at least $2.22$.
\end{abstract}

\tableofcontents
\newpage
\section{Which conclusions use RH?}
Throughout, RH means that every nontrivial zero of the Riemann zeta function
has real part $1/2$. The central distinction is between proving the
implication $\mathrm{RH}\Rightarrow C$ and proving $C$ itself.
Three source questions admit proofs here with no hypothesis on zeta zeros.
\begin{center}
\small\renewcommand{\arraystretch}{1.24}
\begin{tabularx}{\linewidth}{@{}p{.26\linewidth}X@{}}
\toprule
Source question & Conclusion and scope\\
\midrule
Erd\H{o}s--Richards, 1977, Question 3 \cite{er} &
Density one, with all lengths allowed at each good start. The authors
already observed that RH suffices. The proof here removes that assumption.\\
Tom\'as Garc\'ia, 2026, Conjecture 1 \cite{tg} &
The two-sided bound holds for every finite gap multiset. The proposed
lower constant is optimal. The independent small-probability estimate
for the Farey order is quantitative under RH.\\
Csordas, 1998, Problem 4 \cite{cs98} &
A nonreal zero answers the original question in the negative.
Conjugation gives a second zero. This statement concerns the Mellin
transform specified in the question.\\
Csordas, Open Problem 4.11 \cite{cs15} &
For positive deformation parameter, RH implies positive definiteness
at every order. The printed derivative index and the unconditional
eventual assertion are addressed in the proof.\\
Nicolas, 2012, final questions \cite{nic} &
An unconditional improvement for the requested upper bound on the
lower limit. Its exact conjectural endpoint remains undetermined.\\
Chamberland--Straub \cite{chstr} &
A conditional all-start bound at exponent $2.22$. Legendre's
conjecture concerns exponent $2$.\\
\bottomrule
\end{tabularx}
\end{center}
The historical and mathematical questions admit different levels of
verification. Original statements and their quantifiers can be compared
directly with the sources. The proofs and reproducible calculations
below can be checked independently. The sources consulted do not establish
historical priority for these arguments or any ranking over the corpus.

\paragraph{Use of the corpus.}
The repository snapshot \cite{vmrc} contains 225 tracked TeX sources.
We searched this inventory and read the pertinent sections of selected
manuscripts. Item [141] uses exceptional-set estimates for prime counts
on short intervals. Item [149] discusses the large sieve and distinguishes
its additive and multiplicative forms. Both ideas enter the prime-counting
proof below, using published estimates at the required scales.
Item [187] explores failures of iterated Tur\'an tests on multiplier
sequences; the kernel in the present derivative problem is fixed by the
original question. Items [223] and [224] study which positivity transfers
preserve information about zeta zeros. Here an explicit Fourier identity
identifies the tested derivative, while in the Mellin problem a zero-free
normalizing factor preserves the zeros of the original transform.
The repository describes its papers as research drafts. Our theorem
inputs are cited independently of these drafts.

\section{The density conjecture of Erd\H{o}s and Richards}
\label{sec:erdos-richards}
Write $A(n,h)=\pi(n+h)-\pi(n)$. In Question 3 on pp.~110--111 of
\cite{er}, Erd\H{o}s and Richards called an integer $n\ge2$ bad if
$A(n,h)>\pi(h)$ for some integer $2\le h\le n$. They conjectured
that the bad integers have natural density zero and explicitly noted
on p.~111 that RH would suffice. Let $B(x)$ count these integers up to $x$.

\begin{theorem}\label{thm:er-density}
Unconditionally, as $x\to\infty$,
\[
 B(x)\ll x\frac{\log\log x}{\log x}.
\]
Consequently a set of integers $n$ of natural density one satisfies
$A(n,h)\le\pi(h)$ simultaneously for every integer $2\le h\le n$.
\end{theorem}

For comparison, Corollary~1 of Alkan \cite{alkan} estimates the number of
exceptional pairs in a two-dimensional region. A single exceptional
length at each starting point would be compatible with a pair count of
order $X$. Here the exceptional set is projected onto its first
coordinate before counting. The second Hardy--Littlewood conjecture
asks for the inequality at every starting point and every length.
Theorem~\ref{thm:er-density} answers the separate density question of
Erd\H{o}s and Richards.

\subsection{An upper sieve with a mean-square bound}
Fix a large integer $X$ and put $L=\log X$, $R=\lfloor X^{1/4}\rfloor$,
and $D=R^2$. Define
\[
 G=\sum_{r\le R}\frac{\mu^2(r)}{\varphi(r)},\qquad
 y_r=\frac{\mu(r)}{\varphi(r)G}\quad(1\le r\le R),\qquad
 \lambda_d=d\sum_{d\mid r,\ r\le R}\mu(r/d)y_r.
\]
Here $\mu$ is the M\"obius function and $\varphi$ is Euler's totient.
Grouping positive integers by their squarefree radical gives
$G\ge\sum_{m\le R}1/m$, since the total reciprocal mass of all positive
integers with squarefree radical $d$ is $1/\varphi(d)$.
In particular $1/G\le (4+o(1))/L$.

M\"obius inversion gives $\sum_{r\mid d,\ d\le R}\lambda_d/d=y_r$ and
$\lambda_1=1$. The weight vanishes for nonsquarefree $d$. For squarefree
$d$ its absolute value is
\[
 \frac{d}{\varphi(d)G}
   \sum_{m\le R/d,\ (m,d)=1}\frac{\mu^2(m)}{\varphi(m)}\le1.
\]
Indeed the terms of $G$ indexed by $em$, for $e\mid d$ and the squarefree
$m$ in this sum, are distinct, and
$\sum_{e\mid d}1/\varphi(e)=d/\varphi(d)$.
Now set
\[
 W(t)=\Bigl(\sum_{d\mid t,\ d\le R}\lambda_d\Bigr)^2
     =\sum_{\ell\mid t,\ \ell\le D} a_\ell.
\]
This is nonnegative and $W(p)=1$ for every prime $p>R$. The
coefficient $a_\ell$ is the sum of $\lambda_d\lambda_e$ over pairs
with $[d,e]=\ell$, where brackets denote the least common multiple.
It vanishes unless $\ell$ is squarefree, and $|a_\ell|\le\tau_3(\ell)$,
where $\tau_k$ counts ordered factorizations into $k$ positive factors.
There are three assignments per prime divisor of $\ell$ to the two
indices. Its constant Fourier coefficient is
\[
 \sum_{\ell\le D}\frac{a_\ell}{\ell}
 =\sum_{d,e\le R}\frac{\lambda_d\lambda_e}{[d,e]}
 =\sum_{r\le R}\varphi(r)y_r^2=\frac1G.
\]
The middle equality uses $(d,e)=\sum_{r\mid(d,e)}\varphi(r)$.

For every positive integer $H$ the following dispersion estimate holds:
\begin{equation}\label{eq:er-l2}
 \sum_{X\le n<2X}
 \left|\sum_{j=1}^{H}W(n+j)-H/G\right|^2
 \ll X H L^{15}.
\end{equation}
To prove it, expand the divisor indicators by additive characters,
collecting equal fractions in reduced form. If $e(t)=\exp(2\pi it)$
and $K_H(t)=\sum_{j=1}^H e(jt)$, then the coefficient at a reduced
fraction $a/q$, $q\ge2$, of the centered interval sum as a function
of $n$ is $b_q K_H(a/q)$, where
\[
 b_q=\sum_{q\mid\ell,\ \ell\le D}\frac{a_\ell}{\ell},\qquad
 |b_q|\ll \frac{\tau_3(q)L^3}{q}.
\]
The estimate follows from $\tau_3(qm)\le\tau_3(q)\tau_3(m)$ and
$\sum_{m\le D}\tau_3(m)/m\le(1+\log D)^3$. Orthogonality gives
$\sum_{a=1}^{q-1}|K_H(a/q)|^2=r(q-r)\le\min(q^2,qH)$,
where $r$ is the remainder on dividing $H$ by $q$.
The nonzero $b_q$ have squarefree $q$, so that
$\tau_3(q)^2=\tau_9(q)$. Summing the squared coefficients gives
\[
 \ll L^6\left(\sum_{q\le\min(D,H)}\tau_9(q)
             +H\sum_{H<q\le D}\frac{\tau_9(q)}{q}\right)\ll HL^{15}.
\]
The elementary estimates used here follow by summing the nine
factor variables, bounding eight or nine harmonic sums, respectively.
Distinct reduced fractions of denominator at most $D$ are
$D^{-2}$-separated modulo one. The dual additive large sieve
of Montgomery and Vaughan \cite{ls} multiplies the last bound by
at most $X+D^2\le2X$, proving \eqref{eq:er-l2}.

\subsection{All lengths for the same starting point}
Set $H_0=\lceil L^{26}\rceil$. There is an absolute $c>0$ such that
$\pi(h)/h\ge c/\log(2h)$ for all integers $h\ge2$, by the prime
number theorem and adjustment at finitely many integers. If there is
a violation with $2\le h\le H_0$, the average of the prime indicator
on $(n,n+h]$ exceeds $c'/\log L$. This implicates
$O(X\log L/L)$ starting points $n\in[X,2X)$. For a direct covering
argument, take the intervals $[n,n+h]$ of all the witnesses and
select disjoint intervals greedily in order of decreasing length.
Tripling the selected intervals covers their union. On each selected
interval the prime count is $\gg h/\log L$. Their disjointness and
$\pi(3X)\ll X/L$ bound their total length by $O(X\log L/L)$.
Counting integers in their triples changes only the absolute constant.

In any length range $[U,V]$ below we use the integer grid obtained by
taking $H=\lceil U\rceil$ first and then repeatedly replacing $H$ by
$\lceil1.01H\rceil$ until $V$ is reached or crossed. At most $O(L)$
grid points are needed, and for any integer $h\in[U,V]$ the next grid
point satisfies $h\le H\le1.02h$ when $U\ge H_0$.
For $H_0\le h\le X^{1/5}$, apply \eqref{eq:er-l2} on this grid.
The number of starts for which a grid sum exceeds
$H/G+H/(10L)$ is $O(XL^{17}/H)$ at each $H$. Their union has size
$O(XL^{18}/H_0)$. At the remaining starts, monotonicity of prime counts
and the majorant give
\[
 A(n,h)\le \sum_{j=1}^H W(n+j)\le (4+0.1+o(1))H/L
       <4.3h/L<\pi(h).
\]
The last inequality follows uniformly from $\pi(h)\sim h/\log h$,
$h\ge H_0\to\infty$ and $\log h\le L/5$.

For $X^{1/5}\le h\le X^{4/5}$ use a fresh grid.
Guth and Maynard \cite[Corollary~1.4]{gm} give, for each such grid point,
\[
 A(n,H)=H/\log n+O(H\exp(-c_1 L^{1/4}))
\]
except on a set of $O(X\exp(-c_1 L^{1/4}))$ starts, for some $c_1>0$.
Their range is $X^{2/15+\epsilon}\le H\le X^{0.99}$, including the
rounded endpoints for fixed sufficiently small $\epsilon>0$. Their
continuous exceptional-set bound gives the stated integer version:
for integer $H$, both prime counts are constant as the start varies
in $[n,n+1)$, and $H/\log n$ changes by $O(H/(XL^2))$ there.
After taking the union over the grid, all remaining starts satisfy
$A(n,h)\le(1.02+o(1))h/L<\pi(h)$ throughout this range since
$\log h\le4L/5$.

Finally let $X^{4/5}\le h\le n$ and $b=\lfloor X^{7/10}\rfloor$.
Tile $(n,n+h]$ by full blocks of length $b$ and a final segment of
length less than $b$. At every full-block start $u\in[X,4X]$,
\cite[Corollary~1.3]{gm} gives
$\pi(u+b)-\pi(u)=b/\log u+O(b\exp(-c_2 L^{1/4}))$ for a fixed
$c_2>0$, as $u^{17/30+\epsilon}\le b\le u^{0.99}$.
Since $u\ge n$ and $b/h\le X^{-1/10}$, summation yields, uniformly,
\[
 A(n,h)\le h/\log n+o(h/L^2).
\]
The two-term prime number theorem gives
$\pi(h)\ge h/\log h+h/(2\log^2 h)$ for large $h$.
As $h\le n<2X$, this absorbs the error, even when $h=n$.

The dyadic exceptional count is thus $O(X\log L/L)$. Summing over
dyadic intervals between $\sqrt{x}$ and $x$, and counting smaller
starts trivially, proves Theorem~\ref{thm:er-density}. The original
definition starts at $h=2$. If the variant $h=1$ is included, its
extra exceptions have $n+1$ prime and contribute only $O(x/\log x)$.

\paragraph{Literature and scope.}
Erd\H{o}s and Richards give a positive lower density of good
starts and explicitly identify the implication from RH \cite[p.~111]{er}.
We compared their individual-start quantifier with Alkan's pair
count \cite[Corollary~1]{alkan} and checked the short-interval input
against the two separate estimates of Guth and Maynard \cite{gm}.
Selberg's square weights, mean squares for sieve weights in almost
all short intervals, the additive large sieve, and the prime
distribution results used here all have antecedents. For related work
on variance estimates for sieve weights see K.~Matom\"aki,
\emph{Almost primes in almost all very short intervals},
J.~London Math. Soc.\ 106 (2022), 1061--1097,
\href{https://doi.org/10.1112/jlms.12592}{doi:10.1112/jlms.12592},
particularly Sections 1.2 and 5.
The application here uses the numerical overlap between the
sieve bound near $4h/\log X$ and the range of the almost-all prime
number theorem. We have not located an earlier unconditional proof
of the stated individual-start conclusion. Priority would require
further bibliographic and independent mathematical review.

\subsection{How many lengths determine the exceptional set?}
The density theorem leaves a finer question: to a prescribed power of
$1/\log x$, which interval lengths can account for the exceptional starts?
Write $p_j$ for the $j$th prime. Let $B_{\le d}(x)$ count starts up to $x$
with a witness $2\le h\le\min(n,d)$. We obtain a fixed cutoff at every order.
\begin{theorem}\label{thm:projection}
For each fixed integer $K\ge2$,
\[
 0\le B(x)-B_{\le p_K-1}(x)\ll_K \frac{x}{(\log x)^{K+1}}.
\]
\end{theorem}
\begin{proof}
For $X\le n<2X$ write $L=\log X$. In the proof of
Theorem~\ref{thm:er-density}, start the mean-square argument at
$H_0=\lceil L^{K+30}\rceil$. It bounds the exceptional starts with
$h\ge H_0$ by $O_K(XL^{18}/H_0+XL e^{-cL^{1/4}})$.
For $p_K\le h<p_{K+1}$ a violation requires $K+1$ simultaneous primes
with distinct offsets in the fixed interval $[1,p_{K+1}-1]$.
For each admissible collection of $K+1$ offsets, the upper-bound
sieve gives $O_K(X/L^{K+1})$ starts \cite[Lemma 2.3]{lt}.
An inadmissible collection occupies every residue modulo a prime
at most $K+1$ and has no such translates in this dyadic interval for
large $X$. Summing over a fixed finite family preserves the bound.

Put $r=K+2$. Treat $p_{K+1}\le h\le H_0$ in dyadic intervals
$H\le h<2H$ with $H=2^j p_{K+1}$, so that every witness implies
$A(n,2H)\ge r$ and $A(n,2H)\ge \pi(H)+1\gg H/\log(2H)$.
The factorial moment satisfies
\[
 \sum_{X\le n<2X}\binom{A(n,2H)}{r}\ll_{r}\frac{XH^r}{L^r}.
\]
Indeed expand in $r$-element subsets of the offsets. Apply the uniform
upper-bound sieve of \cite[Lemma 2.3]{lt} and sum the singular series
using \cite[Corollary 2.6]{lt}, with the coordinate projections as the
linear forms. Here $2H\ll_K L^{K+30}<X$ for large $X$. Inadmissible
tuples have zero contribution in this interval. For integers $z\ge r$,
$\binom zr\gg_r z^r$. Thus the number of starts with a witness in
one dyadic range is $O_K(X(\log(2H)/L)^r)$. Summing the
$O_K(\log L)$ ranges gives $O_K(X(\log L)^{r+1}/L^r)$.
This is $O_K(X/L^{K+1})$. Summing the estimates over the starting
dyadic intervals proves the theorem.
\end{proof}

Call a finite set of offsets admissible if it omits at least one
residue class modulo every prime. Define
\begin{align*}
 m&=\min\{k\ge2:\ [1,p_k-1]\text{ contains an admissible set of size }k\},\\
 \mathcal F&=\{S\subseteq[1,p_m-1]: |S|=m,\ S\text{ admissible}\},
\end{align*}
where interval subsets in this formula consist of integers.
The existence of $m$ follows from Hensley and Richards
\cite{hensley}: their construction gives an admissible set in $[1,h]$
of cardinality at least $k=\pi(h)+1$ for some $h\ge2$, and $h<p_k$.
Write $T_S(x)=\#\{1\le n\le x : n+s\text{ is prime for every }s\in S\}$.

\begin{corollary}\label{cor:patterns}
Unconditionally,
\[
 B(x)=\sum_{S\in\mathcal F}T_S(x)
       +O_m\!\left(\frac{x}{(\log x)^{m+1}}\right),
 \qquad B(x)\ll_m\frac{x}{(\log x)^m}.
\]
If the usual Hardy--Littlewood asymptotic holds for the fixed tuples in
$\mathcal F$, then $B(x)\sim
(\sum_{S\in\mathcal F}\mathfrak S(S))\,x/(\log x)^m$,
where $\mathfrak S(S)=\prod_p(1-\nu_p(S)/p)(1-1/p)^{-|S|}$
and $\nu_p(S)$ counts occupied residue classes.
\end{corollary}
\begin{proof}
At a sufficiently large $n$, a fixed tuple of $k$ actual primes gives
an admissible set of offsets whenever all these primes exceed $k$.
A witness $h<p_m$ would therefore give, with $k=\pi(h)+1$, an
admissible $k$-set in $[1,h]\subseteq[1,p_k-1]$. Minimality forces
$k=m$. Conversely an occurrence of any $S\in\mathcal F$ is a witness
at $h=p_m-1$ once $n\ge p_m-1$. Two distinct tuples have at least
$m+1$ offsets in their union, giving $O_m(x/(\log x)^{m+1})$
overlaps by the upper-bound sieve. Theorem~\ref{thm:projection} and the
two elementary union bounds give the formula. The upper bound follows
from the same sieve for $m$ offsets. The conditional asymptotic follows
by summing finitely many positive singular series.
\end{proof}
The tuple asymptotics enter only the last assertion. In particular,
the upper bound holds even if no tuple in $\mathcal F$ has infinitely
many prime translates. The implied constants depend on the fixed
order and the result is asymptotic in $x$.

\paragraph{An explicit range for the exponent.}
Let $H(k)$ denote the smallest diameter of an admissible $k$-set,
with $H(1)=0$. Thus a witness of size $k$ is possible exactly when
$H(k)+1<p_k$, for $k\ge2$. The exact diameters for $2\le k\le342$
were determined by Engelsma and are documented in
A.~V.~Sutherland's \emph{Sieve theory and gaps between primes:
Narrow admissible tuples} \cite{suth}, by Polymath \cite[Section 7]{poly}, and in OEIS A008407.
For reproducibility, the table below gives the minimum of
$H(k)+1-p_k$ on successive blocks of the tabulation.
\[
\begin{array}{c|rrrrrrr}
k&2\text{--}50&51\text{--}100&101\text{--}150&151\text{--}200&
201\text{--}250&251\text{--}300&301\text{--}342\\
\hline
\min(H(k)+1-p_k)&0&14&16&28&8&20&4
\end{array}
\]
This is a comparison against published finite computations. To
extend it analytically, a sorted admissible set of size $342+j$
splits into disjoint admissible sets of sizes $342$ and $j$.
Their adjacent endpoints have even positive difference, so
$H(342+j)\ge H(342)+2+H(j)$. With $H(342)=2328$ this gives
\[
\begin{array}{c|rrrrrrrr}
k&343&344&345&346&347&348&349&350\\ \hline
H(k)\ \text{at least}&2330&2332&2336&2338&2342&2346&2350&2356\\
p_k&2309&2311&2333&2339&2341&2347&2351&2357
\end{array}
\]
Meanwhile Engelsma's admissible $447$-tuple of diameter $3158$
gives a witness, since $p_{447}=3163$. Thus
$351\le{m}\le447$. In particular the unconditional estimate is
$B(X)\ll X/(\log X)^{351}$. The data for $k>342$ in tables of
narrowest \emph{known} tuples need not be optimal.
The primary presentation is available at
\url{https://math.mit.edu/~drew/PrimeGapsOberwolfach3.pdf};
the data are at \url{https://oeis.org/A008407/b008407.txt}.

\section{Gap permutations and a conjecture of Tom\'as Garc\'ia}
\label{sec:fareygaps}
Let $0=x_0<x_1<\cdots<x_N=1$, with gaps $g_j=x_j-x_{j-1}$.
The discrepancy considered by Tom\'as Garc\'ia is
\[
 r(S)=\sum_{k=1}^N|x_k-k/N|.
\]
For a uniformly random permutation $\pi$ of the gaps, write $R(S)$
for the expected discrepancy of their successive partial sums. Equal
gaps cause no change in the average: each distinct ordering has the
same number of labelled representations. Put $a_j=g_j-1/N$,
$A=(\sum_j a_j^2)^{1/2}$ and $\sigma_S=A/\sqrt N$.
Conjecture~1 in Section~3 of \cite{tg} asks whether absolute positive
constants bound $R(S)/(\sigma_S N^{3/2})$ from below and above,
uniformly in $S$ with positive variance. The author also proposes two
numerical constants, explicitly designating them provisional.

\begin{theorem}[The permutation estimate]
\label{thm:fareyperm}
For every $N\ge 2$ and every such sequence,
\[
 \frac{N+1}{3\sqrt 2}\,A\ \le\ R(S) \le\sqrt{\frac{N^2-1}{6}}\,A.
\]
The lower coefficient is best possible for each $N$. In the
normalization $\sigma_S N^{3/2}=NA$, the sharp common lower constant
is $1/(3\sqrt2)$. An admissible common upper constant is $1/\sqrt 6$.
\end{theorem}

\noindent
For $N=1$ or $A=0$ the discrepancy and both bounds are zero.
The theorem proves the existential statement in \cite[Conjecture~1]{tg}
and the proposed sharp lower constant. The proposed sharp upper constant
$(4/9)\sqrt{2/3}$ requires a further estimate for general gap multisets.
For comparison, a non-sharp first-moment lower bound for independent
summands is also available from Bobkov and Ledoux \cite[Lemma 3.4]{bl}.

\subsection{Averaging over the size of a random subset}
For $0<p<1$, let $B_1(p),\ldots,B_N(p)$ be independent Bernoulli
variables of success probability $p$. With $T_p=\sum_j a_j(B_j(p)-p)$,
\begin{equation}
 \frac{R(S)}{N+1}=\int_0^1\mathbb E|T_p|\,dp .
 \label{eq:beta-perm}
\end{equation}
Indeed, the first $k$ entries of a uniform permutation form a uniform
$k$-subset. A fixed set $I$ of size $k$ is the set of Bernoulli
successes with probability $p^k(1-p)^{N-k}$. Its integral is
$1/((N+1)\binom Nk)$. Summing over sets gives~\eqref{eq:beta-perm},
since $\sum_j a_j=0$ and the two endpoint partial sums vanish.

\begin{lemma}[A sharp lower estimate at each sampling probability]
\label{lem:bias-bern}
For real $u_1,\ldots,u_N$, independent Bernoulli$(p)$ variables $B_j$
and $0\le p\le1$,
\[
 \mathbb E\left|\sum_j u_j(B_j-p)\right|
 \ge \sqrt2\,p(1-p)\left(\sum_j u_j^2\right)^{1/2}.
\]
Opposite coefficients of equal magnitude, with all other coefficients
zero, give equality for every $p$.
\end{lemma}

\begin{proof}
The endpoints and the zero vector are immediate. Write $q=1-p$ and
normalize $\sum_j u_j^2=1$, omitting zero terms.
If $|u_j|\ge 1/\sqrt2$ for some $j$, conditional Jensen applied
given $B_j$ bounds the left-hand side below by
$|u_j|\mathbb E|B_j-p|=2pq|u_j|$.

Suppose now that $|u_j|\le1/\sqrt2$ for all $j$. Put
$h(v)=1-2pq+2pq\cos v$. This is the characteristic function of the
difference of two independent Bernoulli$(p)$ variables. It is
nonnegative because $2pq\le1/2$. We use the following integral
inequality from Havrilla and Tkocz \cite[proof of Theorem~8,
Section~3.4, equation (23) in the preprint]{ht}. If
$h(v)=\rho_0+(1-\rho_0)\cos v$ with $1/2\le\rho_0\le1$, then
for every real $s\ge1$,
\[
 F_h(s):=\frac2\pi\int_0^\infty
   \bigl(1-h(v/\sqrt s)^s\bigr)\frac{dv}{v^2}
 \ge F_h(1)=1-\rho_0.
\]
Here the atom at zero is $\rho_0=1-2pq$.

Let $T=\sum_j u_j(B_j-p)$, with characteristic function $\phi_T$.
Independence and the weighted arithmetic--geometric mean inequality give
\[
 \operatorname{Re}\phi_T(v)
 \le \prod_j h(u_jv)^{1/2}
 \le \sum_j u_j^2 h(u_jv)^{1/(2u_j^2)} .
\]
The first absolute moment identity
$\mathbb E|T|=(2/\pi)\int_0^\infty(1-\operatorname{Re}\phi_T(v))\,dv/v^2$
follows from the corresponding cosine integral for $|x|$ and Tonelli's
theorem. Apply it to the preceding inequality, use evenness of $h$,
and substitute $v=\sqrt2\,w$. With $s_j=1/(2u_j^2)\ge1$ this yields
\[
 \mathbb E|T|\ge
 \frac1{\sqrt2}\sum_j u_j^2 F_h(s_j)
 \ge \frac{2pq}{\sqrt2}.
\]
The normalized vector $(1,-1,0,\ldots,0)/\sqrt2$ gives
$\mathbb E|B_1-B_2|/\sqrt2=\sqrt2\,pq$.
\end{proof}

\begin{proof}[Proof of Theorem~\ref{thm:fareyperm}]
Integration of Lemma~\ref{lem:bias-bern} in~\eqref{eq:beta-perm}
gives the lower bound, since $\int_0^1p(1-p)\,dp=1/6$.
For its sharpness take $a=(d,-d,0,\ldots,0)$ with $0<d<1/N$.
These centred gaps correspond to a strictly increasing sequence.
Letting $N\to\infty$ gives sharpness of the dimension-independent
constant. For the upper bound, let $S_k=\sum_{j=1}^k a_{\pi(j)}$.
For $i\ne j$ one has $\mathbb E[a_{\pi(i)}a_{\pi(j)}]
=-A^2/(N(N-1))$, so
\[
 \mathbb E S_k^2=\frac{k(N-k)}{N(N-1)}A^2.
\]
Cauchy--Schwarz first in probability and then over $1\le k\le N-1$
gives $R(S)^2\le (N-1)\sum_{k=1}^{N-1}\mathbb E S_k^2
=(N^2-1)A^2/6$.
\end{proof}


\paragraph{The proposed upper constant on a class of gap multisets.}
If just one centered gap is positive, of size $d$, all the negative
magnitudes sum to $d$. Conditional on whether the positive gap was
sampled, the sign of $T_p$ is fixed. The conditional means therefore
give $\mathbb E|T_p|=2p(1-p)d$, so $R=(N+1)d/3$. The same calculation
applies with the signs reversed. By Cauchy--Schwarz,
$d\le A\sqrt{1-1/N}$. The upper ratio in this class is
$(1+1/N)\sqrt{1-1/N}/3$; its maximum for $N\ge2$ occurs at $N=3$,
by differentiation, and equals $(4/9)\sqrt{2/3}$. Choosing all other
centered gaps equal attains it.

\subsection{The leading constant for permuted Farey gaps}
Set $D=\sum_j|a_j|^3$. The same integral formula gives a quantitative
normal approximation for every collection of gaps with $A>0$.
\begin{proposition}
\label{prop:fareynormal}
One has
\[
 \left|\frac{R(S)}{N+1}-\frac{\sqrt{2\pi}}8 A\right|
 \le \frac{2D}{A^2}.
\]
In particular, along any sequence with $N\to\infty$ and
$\max_j|a_j|/A\to0$, the quotient $R(S)/(NA)$ tends to
$\sqrt{2\pi}/8$.
\end{proposition}
\begin{proof}
Write $v_p=p(1-p)A^2$. For fixed $0<p<1$, use the independent-sum
normal approximation for Lipschitz test functions
\cite[Theorem~3.1]{cgs} on $T_p/\sqrt{v_p}$, with test $|x|$.
Multiplying by $\sqrt{v_p}$ gives
\[
 \left|\mathbb E|T_p|-\sqrt{2v_p/\pi}\right|
 \le 3\,\frac{D}{A^2}\bigl(p^2+(1-p)^2\bigr).
\]
Here $\mathbb E|B_j-p|^3=p(1-p)(p^2+(1-p)^2)$.
Use~\eqref{eq:beta-perm} and integrate, with
$\int_0^1\sqrt{p(1-p)}\,dp=\pi/8$.
For the limit use $D/A^3\le\max_j|a_j|/A$.
\end{proof}

Let $F_n$ have its usual positive Farey points through denominator
$n$, with the initial reference point $0/1$. Its number of gaps is
$N=\sum_{q\le n}\varphi(q)=3n^2/\pi^2+O(n\log n)$.
The following computation records the leading term without requiring
a refined asymptotic for the gap moments.
\begin{corollary}
\label{cor:fareyasymp}
Unconditionally, as $n\to\infty$,
\[
 R(F_n)\sim\frac{3\sqrt6}{4\pi^{5/2}}\,n\sqrt{\log n}.
\]
If the ordinary Riemann hypothesis holds, then for every $\varepsilon>0$
\[
 \frac{r(F_n)}{R(F_n)}
  \ll_{\varepsilon}\frac{n^{-1/2+\varepsilon}}{\sqrt{\log n}}.
\]
\end{corollary}
\begin{proof}
Neighbouring fractions of denominators $q,r$ have gap $1/(qr)$.
Their ordered denominator pairs are exactly the coprime pairs
$1\le q,r\le n$ with $q+r>n$. Indeed, the determinant condition
$bq-ar=1$ for fractions $a/q<b/r$ gives a unique pair in $[0,1]$
for each such $(q,r)$, and the smallest possible denominator of
an intervening fraction is $q+r$. Pairs with both denominators $>n/2$
contribute $O(n^{-2})$ to the sum of squared gaps. For $q\le n/2$,
the other denominator runs through $n-q<r\le n$; this interval
contains $\varphi(q)$ values coprime to $q$. Uniformly there
$r^{-2}=n^{-2}+O(q/n^3)$. Counting both orders gives
\[
 \sum_j g_j^2=\frac2{n^2}\sum_{q\le n/2}\frac{\varphi(q)}{q^2}
       + O(n^{-2})
 =\frac{12}{\pi^2}\frac{\log n}{n^2}+O(n^{-2}).
\]
For the last estimate use $\varphi(q)/q=\sum_{d\mid q}\mu(d)/d$,
the harmonic sum, and $\sum_{d\ge1}\mu(d)/d^2=6/\pi^2$.
Consequently $A^2=\sum_jg_j^2-1/N
=(12/\pi^2)n^{-2}\log n+O(n^{-2})$.
The same division of pairs gives
$\sum_jg_j^3\ll n^{-4}+n^{-3}\sum_{q\le n/2}\varphi(q)/q^3
\ll n^{-3}$. Thus $D\ll \sum_j g_j^3+N^{-2}\ll n^{-3}$.
Proposition~\ref{prop:fareynormal} in fact gives
$R(F_n)=(N+1)\sqrt{2\pi}\,A/8+O(n/\log n)$.
Substitution yields the stated leading constant. Finally, Landau's
Farey criterion \cite{landau1924} gives
$r(F_n)\ll_\varepsilon n^{1/2+\varepsilon}$ under RH. Dividing by the
positive unconditional asymptotic proves the last assertion.
\end{proof}

\subsection{The fraction of permutations as regular as the native order}
The mean alone cannot count orderings with small discrepancy. A separate
estimate comes from permuting pairs inside blocks. Let $r_\pi(F_n)$
denote the discrepancy for a uniform ordering of the Farey gaps.

\begin{theorem}[Lower tail for permuted Farey gaps]
\label{thm:fareytail}
There are absolute constants $C,c>0$ such that for $n\ge2$ and $h\ge0$,
\[
 \mathbb P\{r_\pi(F_n)\le h\}
 \le C\exp\!\left[-c\min\{n,n^2/(1+h)^2\}\right].
\]
Consequently, ordinary RH implies that for every $0<\eta<1$ and
all sufficiently large $n$ the fraction of orderings satisfying
$r_\pi(F_n)\le r(F_n)$ is at most $C\exp(-c_\eta n^{1-\eta})$.
\end{theorem}

\begin{proof}
There are two sets of labelled gaps, each of size at least $bN$ for an
absolute $b>0$, such that each value in the larger-gap set exceeds
each value in the smaller-gap set by at least $b/n^2$.
For the smaller gaps take denominator pairs $3n/4<q,r\le n$.
For the larger ones take $n/8<q\le n/4$ and $7n/8<r\le n$.
All these pairs satisfy $q+r>n$. Their gap sizes are respectively at
most $16/(9n^2)$ and at least $4/n^2$. Each rectangle contains
$(6/\pi^2)$ times its area plus $O(n\log n)$ coprime lattice points:
expand coprimality by M\"obius inversion and count the multiples of
each divisor in the two intervals. This proves the assertion for
large $n$; a larger final constant covers bounded $n$.

Divide the first $2Mm$ positions of the permutation into $m$ consecutive
blocks of $2M$ positions, where $M=\lfloor N/(2m)\rfloor$.
Pair the $j$th and $(M+j)$th positions within each block, $1\le j\le M$.
Call a pair good if it contains one gap from each set. Each block has
at least $b^2M$ good pairs except with probability $O(e^{-c_0M})$,
with an absolute $c_0>0$. Indeed a given pair is good with probability
at least $2b^2$. Expose the $2M$ labels of the block successively.
Changing one exposed label while keeping earlier ones fixed admits
a coupling by a transposition of the remaining labels, which alters
the final good-pair count by at most two. Thus the exposure martingale
has increments bounded by two. The Azuma--Hoeffding bound gives
the assertion. The chance that some block fails is $O(me^{-c_0M})$.

Condition on the unordered contents of all pairs and on the remaining
positions. All orientations are now independent fair signs. Put
$S_k=\sum_{j\le k}a_{\pi(j)}$ and let $T_\ell$ denote
the sum of $S_k$ over block $\ell$. Earlier blocks contribute to
$T_\ell$ only via their fixed totals. Thus the $T_\ell$ are conditionally
independent. Swapping the two labels in a pair changes $T_\ell$ by
$M$ times their difference. In the signed-sum expression for $T_\ell$
(up to a fixed shift), every good pair therefore gives a coefficient
of magnitude at least $bM/(2n^2)\gg1/m$, whenever $M\asymp n^2/m$.

If the block has at least $b^2M$ good pairs, its conditional probability
of $|T_\ell|\le 2h/m$ is at most $K(1+h)/\sqrt M$ for an absolute $K$.
Indeed, for $s$ independent fair signs and positive coefficients
at least $d$, the chance of hitting any half-open interval of length $d$
is at most $2^{-s}\binom{s}{\lfloor s/2\rfloor}\ll s^{-1/2}$.
The positive-sign sets for sums in such an interval form an antichain;
this is the Erd\H{o}s--Littlewood--Offord argument \cite{erdoslo}.
Cover the interval of length $4h/m$ by $O(1+h)$ such intervals and
condition on all other signs inside the block.

Take $m=\lfloor\delta\min\{n,n^2/(1+h)^2\}\rfloor$ for a sufficiently
small absolute $\delta>0$. When the quantity inside this floor is
less than two, enlarge $C$ and use probability at most one.
Otherwise $M\asymp n^2/m$, $M\gg m$, and the last block probability
is at most $1/16$, by the choice of $\delta$. Since
$\sum_{\ell=1}^m |T_\ell|\le\sum_{k=1}^N|S_k|$, the event
$r_\pi(F_n)\le h$ forces at least $m/2$ blocks to have
$|T_\ell|\le2h/m$. On conditioning as above and using the union bound,
its probability is at most $2^m(1/16)^{m/2}=2^{-m}$ when every block
has sufficiently many good pairs. Adding $O(me^{-c_0M})$ gives the
first bound. Under RH, Landau's criterion yields
$r(F_n)\ll_\eta n^{1/2+\eta/2}$. Substitute $h=r(F_n)$.
\end{proof}

\section{A nonreal zero of the Mellin transform in Csordas's 1998 problem}\label{sec:mellinzero}

In Problem 4 of \emph{Convexity and the Riemann $\xi$-function}
(1998, p.~49), Csordas asks whether every zero of the meromorphic
continuation of
\[
 M(z)=\int_0^\infty t^{z-1}\Phi(\sqrt t)\,dt,\qquad
 \Phi(t)=\sum_{n=1}^{\infty}\pi n^2(2\pi n^2e^{4t}-3)
                e^{5t-\pi n^2e^{4t}},
 \qquad \Re z>0,
\]
is real and negative. A positive answer would, by a theorem of
P\'olya cited by Csordas, imply RH.

\medskip\noindent
\textbf{Theorem (computer-assisted).}
\emph{The continuation of $M$ has exactly one zero, counted with
multiplicity, in the open disk $|z-z_0|<10^{-26}$, where}
\[
 \begin{split}
 \Re z_0&=-11.6250611666988213619561882819297820434,\\
 \Im z_0&=\phantom{-}6.8381748574844204734881632920867451441 .
 \end{split}
\]
\emph{In particular, there is a simple nonreal conjugate pair of zeros,
and the answer to Csordas's question is negative.}

\medskip\noindent
\emph{Continuation.} Set $R=1/5$ and $F(z)=R^{-2z}M(z)$. The substitution
$t=x^2$ and a split at $x=R$ give
\begin{equation}\label{eq:mellinsplit}
 F(z)=\sum_{k\ge0}\frac{a_k}{z+k}
       +2\int_0^\infty e^{2zu}\Phi(R\,e^u)\,du,\qquad
 a_k=R^{2k}\frac{\Phi^{(2k)}(0)}{(2k)!}.
\end{equation}
Here $R\,e^u$ denotes the product $R\exp(u)$. The series follows from
evenness of $\Phi$ and convergence on $|x|\le R$ of its power
series. The integral is entire by the double-exponential decay of
$\Phi(x)$ for positive $x$. The geometric coefficient bound below
makes the series locally normally convergent off the nonpositive
integers. This proves that \eqref{eq:mellinsplit} gives the required
continuation and allows differentiation under both signs.
The factor $R^{-2z}$ is zero-free.

\medskip\noindent
\emph{Bounds for the series at the origin.}
On $|t|=3/10$ one has $\Re(e^{4t})\ge e^{-1.2}\cos(1.2)>0.1$,
$|e^{4t}|<4$ and $|e^{5t}|<5$. The $n$th summand of $\Phi$ has
modulus bounded by $600 n^4 e^{-0.3 n^2}$. Its sum is below $10^7$.
The sum from $n=29$ onwards is below $10^{-90}$, as follows by
bounding it by twice its first term. Cauchy's estimate then gives
\[
 |a_k|\le10^7(4/9)^k,\qquad |a_k-a_{k,28}|\le10^{-90}(4/9)^k,
\]
where $a_{k,28}$ uses the first 28 theta summands. In the disk of the
theorem $|z+k|>6$ for all $k\ge0$. Replacing the series by
$\sum_{k=0}^{159}a_{k,28}/(z+k)$ costs less than $10^{-48}$ for both
the value and its first derivative. To evaluate the coefficients let
$P_0(X)=2X-3$ and
\[
 P_{j+1}(X)=4X P_j'(X)+(5-4X)P_j(X).
\]
Termwise differentiation at zero gives
$\Phi_{28}^{(j)}(0)=\sum_{n=1}^{28} x_n e^{-x_n}P_j(x_n)$ with
$x_n=\pi n^2$.

\medskip\noindent
\emph{Bounds for the integral.}
For $u\ge2$ put $x=R\,e^u$ and $\eta=e^{4x}>300$. From the defining
series,
\[
 |\Phi(x)|\le 2\pi^2\eta^{9/4}
              \sum_{n\ge1}n^4e^{-\pi n^2\eta}
           \le 1000e^{-\eta}.
\]
For the second bound use $\eta^{9/4}<e^{\eta/2}$ and sum a
geometric majorant. Since $du=d\eta/(\eta\log\eta)$ and
$(2u)^j e^{-22u}<1$ for $u\ge2$, the integral tail and its first
two $z$ derivatives have moduli below $10^{-120}$ throughout the disk.

Partition $[0,2]$ into 32 equal panels. In the complex disk of
radius $1/8$ about the midpoint of any panel, $x=R\,e^u$ satisfies
$|x|<1.68$, $|\Im x|<0.21$, and $\Re x>0$. In particular
$\Re(e^{4x})>0.6$ and $|e^{4x}|<1000$. The theta series gives
$|\Phi(x)|<10^{10}$ there, for example by summing
$10^8 n^4e^{-1.8n^2}$ termwise. In the same disks the integrands
$2(2u)^j e^{2z_0u}\Phi(R\,e^u)$, for $j=0,1$, have modulus
below $10^{15}$: one can use $|2u|<4.3$ and
$|e^{2z_0u}|\le e^{2(12+7)/8}$. On the real panels, replacing $\Phi$
by its first 12 summands costs less than $10^{-100}$ for either
integral. Indeed $x\ge 1/5$ gives $e^{4x}>2$ in the exponential,
while $e^{4x}<400$ and $e^{5x}<2000$ bound the prefactors.

Use Fej\'er's first quadrature rule with $N=102$ nodes on each panel.
Its nodes on $[-1,1]$ are $\cos\theta_j$ with
$\theta_j=(2j+1)\pi/(2N)$, $0\le j<N$. Its weights are
\[
 w_j=\frac2N\left(1-2\sum_{k=1}^{\lfloor(N-1)/2\rfloor}
                   \frac{\cos(2k\theta_j)}{4k^2-1}\right).
\]
Discrete cosine orthogonality and
$\int_{-1}^{1}T_{2k}(v)\,dv=2/(1-4k^2)$ for the Chebyshev
polynomials prove exactness through degree $N-1$. The weights are
positive since $2\sum_{k=1}^{\infty}(4k^2-1)^{-1}=1$, and their
sum equals 2. Each panel has half-width $1/32$, a quarter of the
analytic disk radius. The midpoint Taylor polynomial of degree 101
has uniform error bounded by $10^{15}4^{-102}/(1-1/4)$ on the panel.
Positivity and constant exactness bound the total quadrature error,
summed over all panels, by four times this quantity. The preceding
three truncations and the quadrature error together are below $10^{-44}$
for each of $F(z_0)$ and $F'(z_0)$.

\medskip\noindent
\emph{Arithmetic and zero count.} The interval certificate below uses
110 decimal digits of working precision, encloses $\pi$, and rounds
elementary operations outwards. Let $Q_0,Q_1$ denote the finite
approximations computed there. Its enclosures give
\[
 |Q_0|<2\cdot10^{-40},\qquad \Im Q_1>4.5\cdot10^{-13}.
\]
For reproducibility the real and imaginary midpoints of $Q_0$ are
approximately $1.68841555\cdot10^{-40}$ and
$-1.51587172\cdot10^{-41}$. The analytic error implies
$|F(z_0)|<3\cdot10^{-40}$ and $|F'(z_0)|>4\cdot10^{-13}$.
For $|z-z_0|\le 10^{-26}$, the absolute second derivative of the
series in \eqref{eq:mellinsplit} is at most
$2\cdot10^7/(6^3(1-4/9))$. The central integral contributes at most
$2\cdot2\cdot16\cdot10^{10}$, using $|2u|^2\le16$ and
$|e^{2zu}|\le1$ on $[0,2]$. Together with the far tail this gives
$|F''(z)|<10^{12}$. Hence, on the circle of radius $r=10^{-26}$,
\[
 |F(z)-F'(z_0)(z-z_0)|
 \le 3\cdot10^{-40}+\tfrac12\,10^{12}r^2
 <4\cdot10^{-39}< |F'(z_0)|r .
\]
Rouch\'e's theorem proves the stated count. Schwarz reflection gives
the conjugate zero.\hfill$\square$

\medskip\noindent
\emph{Reading the conclusion.} This refutes the proposed zero location
for the original meromorphic function. The implication in Csordas's
remark uses that zero location as its hypothesis, so its failure gives
no conclusion about RH itself.

\subsection*{Executable interval certificate}
The following Python program uses \texttt{mpmath} 1.3.0. The proof
above bounds integrals and infinite sums analytically. The program
encloses the two finite approximations. Its arithmetic assumptions
are Python's integer operations and the outward enclosures in
\texttt{mpmath.iv}.
{\footnotesize
\begin{verbatim}
from mpmath import iv as I
I.dps = 110
Z = I.mpc
R = I.mpf(1)/5
z = Z('-11.6250611666988213619561882819297820434',
      '6.8381748574844204734881632920867451441')
N = 102
K = 160
L = 32
M = 28
P = [ -3, 2 ]
xs = [I.pi*n*n for n in range(1,M+1)]
wts = [x*I.exp(-x) for x in xs]
f, df = Z(0), Z(0)
def poly(x, p):
    v=I.mpf(p[-1])
    for a in reversed(p[:-1]):
        v = v*x + a
    return v
fact = 1
for k in range(2*K):
    if k%2==0:
        c=R**k/fact*sum(w*poly(x,P) for x,w in zip(xs,wts))
        t=z+k//2
        f += c/t
        df -= c/(t*t)
    P=[(5+4*j)*(P[j] if j<len(P) else 0)
       -4*(P[j-1] if j else 0) for j in range(len(P)+1)]
    fact *= k+1
nodes=[]
for j in range(N):
    theta=I.pi*(2*j+1)/(2*N)
    x=I.cos(theta)
    w=1-2*sum(I.cos(2*k*theta)/(4*k*k-1) for k in range(1,(N-1)//2+1))
    nodes.append((x,2*w/N))
for ell in range(L):
    for x,w in nodes:
        u=(2*ell+1+x)/L
        t=R*I.exp(u)
        eta=I.exp(4*t)
        h=I.exp(5*t)
        phi=h*sum(b*(2*b*eta-3)*I.exp(-b*eta) for b in xs[:12])
        v=2*w/L*I.exp(2*z*u)*phi
        f += v
        df += 2*u*v
print('center approximation',f,flush=True)
print('derivative approximation',df,flush=True)
assert abs(f) < I.mpf('2e-40')
assert df.imag > I.mpf('4.5e-13')
\end{verbatim}
}
\paragraph{Independent arithmetic check.}
A second implementation with 500-bit MPFI/MPFR intervals used 161
Taylor coefficients, 104 quadrature nodes per panel, and 13 summands at each
node. It enclosed the real part of the finite value between
$1.6884155538523815\cdot10^{-40}$ and
$1.6884155538523817\cdot10^{-40}$ and the imaginary part of the
finite derivative between $4.581725629766680\cdot10^{-13}$ and
$4.581725629766681\cdot10^{-13}$. The analytic error bounds above
also cover these increased truncation orders.

\section{Infinitely many strict failures for the Jacobi theta kernel}\label{sec:strict-kernel}
Use the normalization of Csordas's Open Problem 4.13,
\[
 \Phi(t)=\sum_{n\ge1}\pi n^2(2\pi n^2 e^{4t}-3)
                   \exp(5t-\pi n^2e^{4t}),\qquad
 J_m(t)=\Phi^{(m)}(t)^2-\Phi^{(m-1)}(t)\Phi^{(m+1)}(t).
\]
The original Coffey--Csordas conjecture requires $J_m(t)>0$ for every
integer $m\ge1$ and every real $t$. Koide \cite{sra} reports $J_9(0)<0$ and
determines the sign pattern for the ninth order. Sra gives a certified
counterexample at $t=1/50$ in Theorem 5.2 of \emph{GPT, the Counterexample
Machine}. These preprints already disprove the original conjecture.
The following argument explains how strict reversals recur at the same
point for arbitrarily high orders.

\medskip\noindent
\textbf{Theorem.} \emph{All the numbers $\Phi^{(2k)}(0)$, $k\ge0$, are
nonzero. Infinitely many odd integers $m$ satisfy $J_m(0)<0$.
Each of these strict inequalities holds on an open interval about zero
that may depend on $m$.}

\medskip\noindent
\emph{Proof.}
Put $\Theta(x)=\sum_{n\in\mathbb Z}e^{-\pi n^2x}$ for $\Re x>0$, set
$B(t)=e^t\Theta(e^{4t})$, and let $D=d/dt$. Direct termwise differentiation
gives
\[
 (D^2-1)B(t)=16\Phi(t).
\]
Indeed, $(D^2-1)e^{t-u}=8u(2u-3)e^{t-u}$ for $u=\pi n^2 e^{4t}$,
and the terms with indices $n$ and $-n$ coincide. Normal convergence
on compact subsets of $|\Im t|<\pi/8$ justifies differentiation.
Jacobi's identity $\Theta(1/x)=\sqrt{x}\,\Theta(x)$ makes $B$, and hence
$\Phi$, even. Set $a=\pi/8$ and write
\[
 G(y)=\Phi(iy)=\sum_{k\ge0} b_k y^{2k},\qquad
 b_k=\frac{(-1)^k}{(2k)!}\Phi^{(2k)}(0)\qquad (|y|<a).
\]

We first show that $G^{(j)}(y)\to0$ as $y\uparrow a$, for each fixed
$j\ge0$. For $h=a-y$ put
\[
 q=e^{4iy}-i=\sin(4h)+i(\cos(4h)-1),\qquad
 \Re(1/q)=\tfrac12\cot(2h)\sim (4h)^{-1}.
\]
Poisson summation of a Gaussian twisted by $(-1)^n$ gives
\[
 \Theta(i+q)=\sum_{n\in\mathbb Z}(-1)^n e^{-\pi n^2q}
 =q^{-1/2}\sum_{n\in\mathbb Z}
          \exp\!\left(-\frac{\pi(n+1/2)^2}{q}\right)
 \qquad (\Re q>0).
\]
The square root is positive on the positive real axis.
For $h>0$ sufficiently small, $|q|\asymp h$, all derivatives of $q$
are bounded, and $\Re(1/q)\ge c/h$ for a constant $c>0$.
At any fixed differentiation order the absolute values of the terms
after differentiation with respect to $y$ are bounded by
\[
 C h^{-A}(1+|n|)^A \exp\bigl(-\pi c(n+1/2)^2/h\bigr)
\]
for some $A,C$. The sum is $O(h^{-A}e^{-c'/h})$ for some $c'>0$, by
splitting the exponential in half and summing the remaining Gaussian
at a fixed upper bound on $h$. This proves flatness of $B(iy)$ and all
its derivatives. The differential identity proves the assertion for $G$.

We next exclude zero coefficients. Romik's centered theta expansion \cite{romik}
(Theorem 2 and equations (7)--(8) of \emph{The Taylor coefficients of
the Jacobi theta constant $\theta_3$}), with $x=e^{4t}$, gives
\[
 \frac{B(t)}{\Theta(1)}=
 (\cosh 2t)^{-1/2}
 \sum_{j\ge0}\frac{d(j)}{(2j)!}\,\kappa^j(\tanh 2t)^{2j},\qquad
 \kappa=\frac{\Gamma(1/4)^8}{128\pi^4},
\]
locally at zero, where $d(0)=1$ and all $d(j)$ are integers. Thus
$16\Phi^{(2k)}(0)/\Theta(1)=P_k(\kappa)$ for a polynomial
$P_k\in\mathbb Q[X]$. Only $j\le k+1$ contributes. For
$g(t)=(\cosh 2t)^{-1/2}$ its constant term is
\[
 P_k(0)=g^{(2k+2)}(0)-g^{(2k)}(0).
\]
To see that it is nonzero, write $A(y)=g(iy)=(\cos 2y)^{-1/2}$.
All its even Taylor coefficients are strictly positive. Indeed,
$T(y)=\tan(2y)$ has strictly positive odd coefficients by
$T'=2(1+T^2)$ and $T(0)=0$, and then $A'=TA$, $A(0)=1$ gives the claim
by induction. Consequently $P_k(0)=(-1)^{k+1}
(A^{(2k+2)}(0)+A^{(2k)}(0))\ne0$. Chudnovsky's algebraic
independence \cite[p.~399]{roy} of $\pi$ and $\Gamma(1/4)$ implies that $\kappa$ is
transcendental. It follows that $P_k(\kappa)\ne0$ for every $k$.

Suppose the nonzero real coefficients $b_k$ eventually have one sign.
Choose $K$ so that they have a common sign for all $k\ge K$.
Termwise differentiation of the convergent Taylor series gives
$|G^{(2K)}(y)|\ge (2K)!|b_K|$ for $0<y<a$. Flatness gives a
contradiction. Thus $b_k$ takes both signs infinitely often. Since
no $b_k$ vanishes, $b_kb_{k+1}<0$ infinitely often. Evenness yields
\[
 J_{2k+1}(0)=-\Phi^{(2k)}(0)\Phi^{(2k+2)}(0)
 =(2k)!(2k+2)!\, b_kb_{k+1}.
\]
Continuity gives the asserted neighborhoods.\hfill$\square$

\medskip\noindent
The theorem uses no hypothesis about zeta zeros. The inequalities involve
derivatives of the Fourier kernel $\Phi$. Csordas's distinct Open Problem
4.7 concerns derivatives of its Fourier transform and remains outside this
argument. The prior finite-order counterexamples rule out a first
resolution of Problem 4.13. Priority for the infinite-order statement
has not been established.
In particular the residue of $M(z)$ at $z=-k$ is
$\Phi^{(2k)}(0)/(2k)!$ by \eqref{eq:mellinsplit}. Its nonvanishing shows
that every candidate pole of that continuation is genuine.

\section{The kernel in Csordas's Open Problem 4.11}
Let $\Phi$ be the even Jacobi kernel in Csordas's normalization. Write
$\phi(t)=e^{\lambda t^2}\Phi(t)$ and
$H_\lambda(x)=\int_0^\infty\phi(t)\cos(xt)\,dt$. In Open Problem~4.11
of \emph{Fourier transforms of positive definite kernels and the Riemann
$\xi$-function}, Csordas \cite{cs15} asks about the positive definiteness of the
following specific kernel for $\lambda>0$:
\[
 C_{\lambda,m}(t)=\int_{\mathbb R}\phi(s+t)\phi(s-t)
                        (s^2-t^2)^m s^2\,ds.
\]
We use the exponent $m$ printed in the question. For every $\lambda>0$
there is a threshold $M(\lambda)$ beyond which the kernel is positive
definite without any hypothesis. Assuming RH, the conclusion holds for
every $\lambda>0$ and every integer $m\ge0$, including all orders
covered by the question's preceding assumptions.

\paragraph{Proof.} Use $\widehat f(\omega)=\int_{\mathbb R}f(t)e^{i\omega t}\,dt$.
The Jacobi kernel and its derivatives decrease super-exponentially, so all the
weighted functions below are Schwartz functions. Changing variables to
$u=s+t$, $v=s-t$ gives $ds\,dt=du\,dv/2$, $s^2=(u+v)^2/4$, and
$s^2-t^2=uv$. If $I_j(x)=\int_{\mathbb R}t^j\phi(t)e^{ixt}\,dt$, then
$I_j(x)=2i^{-j}H_\lambda^{(j)}(x)$ and $I_j(-x)=(-1)^jI_j(x)$.
Absolute convergence permits us to expand $(u+v)^2$ before integration.
Thus
\begin{align*}
 8\widehat C_{\lambda,m}(2x)
 &= I_{m+2}(x)I_m(-x)+2I_{m+1}(x)I_{m+1}(-x)+I_m(x)I_{m+2}(-x),\\
 \widehat C_{\lambda,m}(2x)
 &= (H_\lambda^{(m+1)}(x))^2-H_\lambda^{(m)}(x)H_\lambda^{(m+2)}(x).
\end{align*}
For every real entire function $f$ in the Laguerre--P\'olya class,
$(f')^2-ff''\ge0$ on the real axis. This follows first for real-rooted
polynomials by their logarithmic derivatives, with the values at zeros
obtained by continuity, and then by locally uniform approximation.
Fourier inversion therefore gives positive definiteness of
$C_{\lambda,m}$ whenever $H_\lambda^{(m)}$ is in this class.
Indeed, its quadratic form at any finite set of points is the integral
of a squared exponential sum against the nonnegative density
$\widehat C_{\lambda,m}(\omega)/(2\pi)$.

For fixed $\lambda>0$, Ki, Kim and Lee show that $H_\lambda$ has only
finitely many nonreal zeros. The P\'olya--Wiman theorem gives
$H_\lambda^{(m)}$ in the Laguerre--P\'olya class for every $m$ sufficiently
large, as recalled just above (4.12) in Csordas. Under RH, $H_0$ belongs
to the class and de Bruijn's real-zero preservation places $H_\lambda$
in it for every $\lambda\ge0$. Differentiation preserves the class,
proving the all-orders assertion.\hfill$\square$

The derivative index matters. Equation (4.12) immediately before the
question concerns order $2m$. The kernel as printed tests order $m$.
The unconditional argument works for every $m\ge M(\lambda)$. It also
answers the eventual-order reading if the existential threshold in the
preceding paragraph may be enlarged to $2M(\lambda)$. From the single
fact that $H_\lambda^{(2m)}$ belongs to the class we do not deduce the
same assertion at order $m$. Under RH the distinction disappears.
The transform identity specializes the correlation principle behind
Csordas's Theorem~3.7; related general Wronskian transform identities
are due to Dimitrov and Xu (\emph{Trans. Amer. Math. Soc.},
\href{https://doi.org/10.1090/tran/7809}{doi:10.1090/tran/7809}).
The argument establishes the implication, without a claim of a previously
unknown RH method or priority for this specialization.

Two adjacent questions have different histories. In Remark~4.8 Csordas
already identifies RH as sufficient for Open Problem~4.7, the inequality
$(H'_0)^2-H_0H''_0\ge0$ on $\mathbb R$. For Open Problem~4.10 the
Rodgers--Tao theorem \cite{rt} $\Lambda\ge0$ gives a nonreal zero of $H_\lambda$
whenever $\lambda<0$. Here $e^{\lambda t^2}\Phi(t)$ retains the strict
log concavity on $t>0$ and admissibility required in Csordas's
Theorem~3.7. Its criterion therefore forces at least one of the associated
kernels of (4.11) in that paper to fail positive definiteness. This last
existence conclusion is unconditional and follows from established results.

\section{An unconditional bound for the oscillation asked about by Nicolas}
Write $\rho=\sigma+i\tau$ for a nontrivial zero, counted with multiplicity,
and $\gamma_E$ for Euler's constant. Nicolas (2012, (1.18) and p.~320)
defined
\[
 W(x)=\sum_\rho\frac{x^{i\operatorname{Im}\rho}}{\rho(1-\rho)},
 \qquad B=2+\gamma_E-\log(4\pi)=0.0461914179\ldots
\]
and asked for a smaller upper bound for $\liminf W$ than
$-1/(1/4+\tau_1^2)=-0.00499\ldots$. He separately asked whether the
exact value is $-B$. The following answers his quantitative request.

\paragraph{Proposition.}
Without assuming RH,
\[
 -B\ \le\ \liminf_{x\to\infty}W(x)\ <\ -0.0431853 .
\]
For $P_k=\prod_{j\le k}p_j$, put
$c(P_k)=(P_k/\varphi(P_k)-e^{\gamma_E}\log\log P_k)\sqrt{\log P_k}$.
Then, unconditionally, $\liminf_k c(P_k)<3.485229$.

\paragraph{Proof.}
Set $F(v)=W(e^v)$. The zero-counting formula gives absolute and uniform
convergence. Functional-equation symmetry pairs $\sigma+i\tau$ with
$1-\sigma+i\tau$. Thus $F$ is a uniformly almost periodic cosine series
with coefficients represented by the positive measure
\[
 d\mu(t)=2\!\sum_{\operatorname{Im}\rho>0}
 \operatorname{Re}\frac1{\rho(1-\rho)}\,\delta_{\operatorname{Im}\rho}(dt),
 \qquad F(v)=\int\cos(vt)\,d\mu(t).
\]
Indeed, writing $a=t^2+\sigma(1-\sigma)$ and $b=t(1-2\sigma)$,
the real part of the reciprocal is $1/(a+b^2/a)$ and lies between
$1/(t^2+5/4)$ and $1/t^2$. The classical Hadamard-product evaluation
$\sum_\rho 1/(\rho(1-\rho))=B$ gives $\mu((0,\infty))=B$.
In particular, $|F|\le B$ without RH. Uniform almost periodicity gives
$\liminf_{v\to\infty}F(v)=\inf_{v\in\mathbb R}F(v)$:
apply simultaneous homogeneous approximation to any finite set of
frequencies, then let the uniform tail tend to zero.

The inhomogeneous approximation of Mossinghoff and Trudgian
(\emph{Oscillations in the Goldbach conjecture}, (3.1), Table~4.1 and
Fig.~4.1, DOI \texttt{10.5802/jtnb.1202}) supplies a real $y$ for which
all the first 650 positive ordinates have $y\tau_j$ within
$\epsilon=0.014$ of odd multiples of $\pi$. We have rounded the reported phase error upwards. This is finite information.
The $650$th ordinate exceeds $1000$ ($\tau_{650}=1001.34948\ldots$;
Platt's zero data). Hence $\cos(yt)\le-1+\epsilon^2/2$
on the support up to $T=1000$.

Average $F(y+u)$ against the centered Gaussian of variance
$2s$, with $s=3\cdot10^{-7}$. Its multiplier is $h(t)=e^{-st^2}$.
The global infimum of $F$ is at most that average. For the mass up to $T$,
$1-h(t)\le st^2$ and each positive-ordinate summand of $t^2\,d\mu$
has mass at most $2$. Writing $A(t)=\mu((t,\infty))$, this bounds the
low-frequency contribution by
\[
 -(1-\epsilon^2/2)B + A(T)+1300s .
\]
For $T_j=1000+100j$, $0\le j\le70$, the high-frequency contribution is
at most
\[
 h(T)A(T)-
 \sum_{j=0}^{69}\{h(T_j)-h(T_{j+1})\}A(T_{j+1}).
\]
This follows by majorizing $h$ by its left-endpoint values on the
successive intervals, with constant value $h(T_{70})$ afterwards.

For $t\ge1000$ put
\[
 \begin{aligned} b(t)&=\frac{1+\log(t/(2\pi))}{2\pi t},\qquad
 e(t)&=\frac{1+4\log t}{t^2},\\
 L(t)&=\frac{2(b(t)-e(t))}{1+5/(4t^2)},\qquad U(t)&=2(b(t)+e(t)).\end{aligned}
\]
Lehman's unconditional zero-counting lemma, applied to $u^{-2}$, gives
$|\sum_{\tau>t}\tau^{-2}-b(t)|\le e(t)$
(also Mossinghoff--Trudgian, (3.4)). The preceding bounds on each
real reciprocal give $L(t)\le A(t)\le U(t)$. The resulting upper bound is
\[
 D=-(1-\epsilon^2/2)B+(1+h(T))U(T)+1300s-
 \sum_{j=0}^{69}\{h(T_j)-h(T_{j+1})\}L(T_{j+1}).
\]
Direct interval evaluation gives
$-0.043185313913<D<-0.043185313910$.
All constants after the finite phase certificate use elementary functions.
The following Python calculation with \texttt{mpmath.iv} reproduces the enclosure.
\begin{verbatim}
from mpmath import iv as I
I.dps = 50
p = I.pi
B = 2+I.euler-I.log(4*p)
s, ep = I.mpf(3)/10**7, I.mpf(14)/1000
def h(t): return I.exp(-s*t**2)
def b(t): return (1+I.log(I.mpf(t)/(2*p)))/(2*p*t)
def e(t): return (1+4*I.log(t))/t**2
def lo(t): return 2*(b(t)-e(t))/(1+I.mpf(5)/(4*t**2))
def hi(t): return 2*(b(t)+e(t))
D = -(1-ep**2/2)*B+(1+h(1000))*hi(1000)+1300*s
for t in range(1000,8000,100):
    D -= (h(t)-h(t+100))*lo(t+100)
assert I.mpf('-0.043185313913') < D < I.mpf('-0.043185313910')
\end{verbatim}
Under RH, Nicolas (2012, Proposition~3.1) gives
$c(P_k)=e^{\gamma_E}(2+W(p_k))+o(1)$. Uniform continuity of $F$ and
$\log(p_{k+1}/p_k)\to0$ show that the liminf of $W(p_k)$ equals that of
$W(x)$ as $x\to\infty$. If RH fails, Nicolas's oscillation result (2012,
(1.10) and the calculation immediately after it) instead gives
$\liminf_k c(P_k)=-\infty$: $f(x)$ is constant between consecutive primes,
$\log f(x)=\Omega_{\pm}(x^{-b})$ for some $b<1/2$, and
$c(P_k)\sim-e^{\gamma_E}\sqrt{p_k}\log p_k\log f(p_k)$ along those
oscillating sequences. Both cases give the stated strict upper bound
since $e^{\gamma_E}(2+D)<3.485229$.\hfill$\square$

\paragraph{Scope and provenance.}
The absolute value of the bound is more than $93\%$ of $B$.
The inhomogeneous phase certificate is the published one of
Mossinghoff and Trudgian. Without any Gaussian average the same data
already give a bound below $-0.04220$ by $-(1-\epsilon^2/2)B+2U(1000)$.
The improvement above uses the spectral mass between successive
thresholds after Gaussian averaging. Neither argument forces
$\liminf W=-B$. An independent resolution of the request for a smaller
bound, or an earlier numerical value in other notation, remains possible.
\paragraph{A full interval of subsequential limits.}
Every number in $[3.485229,e^{\gamma_E}(2+B)]$ is, unconditionally,
a subsequential limit of $c(P_k)$. Under RH the endpoint on the right
is the primorial limsup by Nicolas \cite[Section 4]{nic}. If RH fails
the corresponding limsup is $+\infty$ and the liminf is $-\infty$,
by the two signs in Nicolas's oscillation theorem. In either case
the consecutive increments of $c(P_k)$ tend to zero. To see this, write
$\theta=\log P_k\asymp p_k$, $R_k=P_k/\varphi(P_k)\ll\log p_k$,
and $q=p_{k+1}\sim p_k$. The changes in $\theta$ and $R_k$ are
$\log q$ and $R_k/(q-1)$. The mean-value theorem applied to
$R_k\sqrt{\theta}-e^{\gamma_E}\sqrt{\theta}\log\theta$ bounds the
increment by $O((\log p_k)^2/\sqrt{p_k})$. A sequence whose increments
tend to zero visits every interval separating its liminf from its
limsup infinitely often, and finite endpoints are subsequential limits.
This proves the claim.

\section{A conditional bound for primes between consecutive powers}
Chamberland and Straub \cite{chstr} proved that RH gives a prime in every interval
$[x^A,(x+1)^A]$ for real $x\ge1$ and $A\ge 9/4$. At the end
of their Section~3 they recorded the sharper threshold $2.2253$.
Retaining a logarithm in the prime-gap estimate gives a further reduction.

\paragraph{Proposition.} Assume RH for the Riemann zeta function. For every real
$A\ge 111/50$ and $x\ge1$, the interval $[x^A,(x+1)^A]$ contains a prime.

\paragraph{Proof.} First take $A=111/50$. Put $Y=x^A$ and
$N=7.05\cdot10^{13}$. For $1\le\sqrt Y<N$, Proposition~3.4 of Chamberland and
Straub, using the computation of Sorenson and Webster, supplies a prime in
$[Y,(\sqrt Y+1)^2]$. Since $(x+1)^{A/2}\ge x^{A/2}+1$, this prime belongs
to the required interval.

Suppose $Y\ge N^2$. Let $L=\log Y$, $L_0=2\log N$, $b=11/222$,
$G=(A/2)e^{bL}$, and $c=2G/L$. The interval from $Y$ to
$Y+c\sqrt Y\,L=Y+Ax^{A-1}$ is contained in $[x^A,(x+1)^A]$ by convexity.
Since $bL>1$, $c$ increases with $L$, and $c(L_0)>0.81$. For $c\ge1$ the
conclusion follows from Theorem~5 of Carneiro, Milinovich and Soundararajan,
which gives a prime in $[Y,Y+\frac{22}{25}\sqrt Y\log Y]$.

For $1/2<c\le1$ suppose that $[Y,Y+c\sqrt Y\,L]$ contains no prime.
Put $u=cLe^{-L/2}$, $a=Y(1+u)^{1/2}$, and $d=\tfrac12\log(1+u)$.
Equations (5.9)--(5.11), (6.11), and (6.16) in the authors' preprint \cite{cms}
\emph{Fourier optimization and prime gaps}, together with the displayed
inequality immediately following (6.16), give
\[
 J\sqrt a\,d\le \log(1/d)+1.16,\qquad J=1.1405\ldots>1.14,
 \qquad \sqrt a\,d\ge G-0.001.
\]
Their assumptions are RH, $Y\ge4\cdot10^{18}$, $1/2\le c\le1$,
and absence of a prime in this closed interval. Since
$u\le Le^{-L/2}<10^{-11}$, the bound
$\log(1+u)\ge u(1-u/2)$ gives $\log(1/d)\le L/2-\log G+0.001$.
Consequently $h(L)=1.14(G-0.001)+\log G-L/2-1.161\le0$.
But $63.77<L_0<63.78$, $G(L_0)>26.15$, and $\log G(L_0)>3.263$ imply
$h(L_0)>0.02$. Moreover $h'(L)=1.14bG+b-1/2>0$ for $L\ge L_0$.
This contradiction completes the proof for $A=111/50$.
For $B>A$, take $t=x^{B/A}$ and use $t+1\le(x+1)^{B/A}$ to include
$[t^A,(t+1)^A]$ in $[x^B,(x+1)^B]$.\hfill$\square$

The comparison concerns the explicit exponent recorded by Chamberland and
Straub. The conjecture for squares, $A=2$ and integer $x$, remains open.

\section{Scope of the implications and prior results}
A true conclusion proved unconditionally also follows under RH. This
logical observation alone supplies no historical novelty. For the
Erd\H{o}s--Richards question the conditional direction appears in the
1977 paper itself; the gain in Section~\ref{sec:erdos-richards} is the
removal of RH and a reduction of all exceptional lengths to finitely
many fixed constellations at each logarithmic order. The explicit
exponent uses independently published finite data. The main density
proof and the finite-length reduction do not require those data.

For Tom\'as Garc\'ia's quantitative conjecture the universal existence
of the constants and the optimal lower constant are proved without RH.
The sharp upper constant suggested by experiments remains undecided
here for arbitrary multisets. A lower comparison of the mean does not
control the number of nearly as regular permutations; the independent
pair orientations in Theorem~\ref{thm:fareytail} supply that control.
For Csordas's 1998 question the proof gives a negative answer in the
precise meromorphic continuation specified by the author. Its
sufficient condition for RH therefore fails irrespective of the
truth of RH itself.

Several direct comparisons also rule out priority. Saito's revised
Theorem~1.8 \cite{saito} already deduces transcendence of the original
Mills constant from the density hypothesis for the zeta function,
which follows from RH. Rodgers and Tao \cite{rt} proved
$\Lambda\ge0$ unconditionally, while the classical real-zero threshold
gives $\mathrm{RH}\Leftrightarrow\Lambda\le0$. Hence the implication
$\mathrm{RH}\Rightarrow\Lambda=0$ also has prior proofs. Neither
serves as a new conditional discovery here.
For a neighboring kernel problem, Rainstar publicly archived a proof
of strict total positivity of order four for the translation kernel
in July 2026 \cite{pf4}. That normalization is $2\Phi(t/2)$.
Reversing the second set of nodes gives the corresponding sign
regularity for $\Phi(u+v)$; we do not count it as a first solution.
The Mellin transform in Csordas's 1998 Problem 4 is a separate function.

The six criteria in the research question separate source identity,
preservation of the original statement, previous open status, unconditional
proof, mathematical gain and priority. Sections~\ref{sec:erdos-richards},
\ref{sec:fareygaps} and \ref{sec:mellinzero} return to the statements
in the cited primary sources. Their proposed solutions are open to
external review. Searches of the source titles, formulations, adjacent
problems and subsequent work did not identify an earlier proof of these
three conclusions with the same quantifiers. This establishes the
scope of the comparison; it cannot certify first discovery or
establish a new historical principle. The all-start, all-length
Hardy--Littlewood subadditivity problem and RH remain open in this
manuscript.

\begin{thebibliography}{99}
\bibitem{er} P. Erd\H{o}s and I. Richards, Density functions for prime and relatively prime numbers, \emph{Monatsh. Math.} \textbf{83} (1977), 99--112. \href{https://renyi.hu/~p_erdos/1977-09.pdf}{Original paper}, Question 3.
\bibitem{alkan} E. Alkan, A generalization of the Hardy--Littlewood conjecture, \emph{Integers} \textbf{22} (2022), A53. \href{https://math.colgate.edu/~integers/w53/w53.pdf}{Article}.
\bibitem{ls} H. L. Montgomery and R. C. Vaughan, The large sieve, \emph{Mathematika} \textbf{20} (1973), 119--134. \href{https://doi.org/10.1112/S0025579300004708}{DOI}.
\bibitem{gm} L. Guth and J. Maynard, New large value estimates for Dirichlet polynomials, \emph{Ann. of Math.} \textbf{203} (2026), 623--675. \href{https://doi.org/10.4007/annals.2026.203.2.6}{DOI}; \href{https://arxiv.org/html/2405.20552v1}{preprint}, Corollaries 1.3 and 1.4.
\bibitem{lt} J. D. Lichtman and J. Ter\"av\"ainen, On the Hardy--Littlewood--Chowla conjecture on average, \emph{Forum Math. Sigma} \textbf{10} (2022). \href{https://doi.org/10.1017/fms.2022.54}{DOI}, Lemma 2.3 and Corollary 2.6.
\bibitem{hensley} D. Hensley and I. Richards, Primes in intervals, \emph{Acta Arith.} \textbf{25} (1974), 375--391.
\bibitem{poly} D. H. J. Polymath, Variants of the Selberg sieve, and bounded intervals containing many primes, \emph{Res. Math. Sci.} \textbf{1} (2014), article 12. \href{https://doi.org/10.1186/s40687-014-0012-7}{DOI}.
\bibitem{suth} A. V. Sutherland, \emph{Sieve theory and gaps between primes: Narrow admissible tuples}, Oberwolfach lecture, 2015. \href{https://math.mit.edu/~drew/PrimeGapsOberwolfach3.pdf}{Slides}. T. J. Engelsma, \emph{Permissible patterns of primes}, 2009 working manuscript, Table 5. \href{https://www.opertech.com/primes/permissiblepatterns.pdf}{Table}.
\bibitem{clark} D. A. Clark and N. C. Jarvis, Dense admissible sequences, \emph{Math. Comp.} \textbf{70} (2001), 1713--1718. \href{https://doi.org/10.1090/S0025-5718-01-01348-5}{DOI}.
\bibitem{cs98} G. Csordas, Convexity and the Riemann $\xi$-function,
\emph{Glas. Mat.} \textbf{33(53)} (1998), 37--50. \href{https://web.math.pmf.unizg.hr/glasnik/33.1/33105.pdf}{Article}, Problem 4, p.~49.
\bibitem{cs15} G. Csordas, Fourier transforms of positive definite kernels and the Riemann $\xi$-function, \emph{Comput. Methods Funct. Theory} \textbf{15} (2015), 373--391. \href{https://arxiv.org/abs/1309.0055}{Author's preprint}.
\bibitem{romik} D. Romik, The Taylor coefficients of the Jacobi theta constant $\theta_3$. \href{https://arxiv.org/abs/1807.06130}{arXiv:1807.06130}, Theorem 2.
\bibitem{roy} D. Roy and M. Waldschmidt, Simultaneous approximation and algebraic independence, \emph{Ramanujan J.} \textbf{1} (1997), 379--430. \href{https://webusers.imj-prg.fr/~michel.waldschmidt/articles/pdf/RoyWaldschmidtRamaJ1.pdf}{Article}, p.~399, citing G. V. Chudnovsky.
\bibitem{sra} S. Sra, \emph{GPT, the Counterexample Machine}, 2026 preprint, \href{https://arxiv.org/html/2608.29595v1}{Theorem 5.2}. A. Koide, independent ninth-order analysis, \href{https://doi.org/10.5281/zenodo.21944522}{Zenodo preprint}.
\bibitem{rt} B. Rodgers and T. Tao, The de Bruijn--Newman constant is non-negative, \emph{Forum Math. Pi} \textbf{8} (2020), e6. \href{https://doi.org/10.1017/fmp.2020.6}{DOI}.
\bibitem{nic} J.-L. Nicolas, Small values of the Euler function and the Riemann hypothesis, \emph{Acta Arith.} \textbf{155} (2012), 311--321. \href{https://arxiv.org/pdf/1202.0729}{Preprint}.
\bibitem{mt} M. J. Mossinghoff and T. S. Trudgian, Oscillations in the Goldbach conjecture, \emph{J. Th\'eor. Nombres Bordeaux} \textbf{34} (2022), 295--307. \href{https://doi.org/10.5802/jtnb.1202}{DOI}.
\bibitem{chstr} M. Chamberland and A. Straub, Weakening the Legendre conjecture, \emph{Amer. Math. Monthly} (2026). \href{https://doi.org/10.1080/00029890.2026.2694309}{DOI}; \href{https://arxiv.org/abs/2602.22502}{preprint}.
\bibitem{cms} E. Carneiro, M. B. Milinovich and K. Soundararajan, Fourier optimization and prime gaps, \emph{Comment. Math. Helv.} \textbf{94} (2019), 533--568. \href{https://arxiv.org/abs/1708.04122}{Preprint}.
\bibitem{saito} K. Saito, Transcendency of variants of Mills' constant, \href{https://arxiv.org/html/2508.16068v3}{arXiv:2508.16068v3}, Theorem 1.8.
\bibitem{vmrc} QEEDDD, \emph{Verifiable Mathematical Reasoning Corpus},
\href{https://github.com/QEEDDD/Verifiable-Mathematical-Reasoning-Corpus/tree/c2ac34ea4622778187173f41bb1996deee8081b7}{fixed revision c2ac34ea4622}, 20 November 2026.


\bibitem{pf4} J. Rainstar, \emph{The Riemann kernel is globally PF4}, revised 22 July 2026, \href{https://doi.org/10.5281/zenodo.21487371}{archived paper and certificates}.
\bibitem{bl} S. Bobkov and M. Ledoux, One-dimensional empirical measures, order statistics, and Kantorovich transport distances, \emph{Mem. Amer. Math. Soc.} \textbf{261} (2019), no.~1259. \href{https://doi.org/10.1090/memo/1259}{DOI}.

\bibitem{tg}
R. Tom\'as Garc\'ia, A general lower bound for average local discrepancy and an application to the Farey sequence, \emph{Mathematics} \textbf{14} (2026), 2543. DOI 10.3390/math14142543.
\bibitem{ht}
A. Havrilla and T. Tkocz, Sharp Khinchin-type inequalities for symmetric discrete uniform random variables, \emph{Israel J. Math.} \textbf{246} (2021), 281--297. DOI 10.1007/s11856-021-2244-8. Preprint arXiv:1912.13345.
\bibitem{cgs}
L. H. Y. Chen, L. Goldstein and Q.-M. Shao, \emph{Normal Approximation by Stein's Method}, Springer, Heidelberg, 2011. DOI 10.1007/978-3-642-15007-4.
\bibitem{landau1924}
E. Landau, Bemerkungen zu der vorstehenden Abhandlung von Herrn Franel, \emph{Nachr. Ges. Wiss. G\"ottingen, Math.-Phys. Kl.} (1924), 202--206.

\bibitem{erdoslo}
P. Erd\H{o}s, On a lemma of Littlewood and Offord, \emph{Bull. Amer. Math. Soc.} \textbf{51} (1945), 898--902.

\end{thebibliography}
\end{document}
